\documentclass[11pt]{article}

\usepackage[margin=1in]{geometry}
\usepackage{amsmath,amssymb}

\title{A Question About a Cyclic Ordering of Type \(A\) Chambers}
\author{Yehonathan Sharvit}
\date{}

\begin{document}
\maketitle

Dear Prof. Baruch,

I would like to ask whether the following formulation sounds natural or
interesting from the viewpoint of Weyl groups.

Consider the braid arrangement
\[
  \mathcal A_{n-1}=\{x_i=x_j\}
\]
in
\[
  \mathbb R^n/\langle(1,\ldots,1)\rangle.
\]
Its chambers are indexed by the \(n!\) possible linear orders of the
coordinates \(x_1,\ldots,x_n\).

Using the known recursive Hamiltonian-cycle algorithm of Zaks for permutations,
one obtains a cyclic ordering of all these \(n!\) chambers. What seems striking
is that this ordering is equivariant for the natural dihedral action on the
indices \(1,\ldots,n\): the cyclic rotation of the indices acts as a rotation of
the large cycle of chambers, and the reversal of the indices acts as a
reflection.

Equivalently, in type \(A_{n-1}\), the ordering is compatible with the natural
dihedral group generated by a Coxeter element and the longest element \(w_0\).

My question is simply: is such a recursive, dihedrally equivariant cyclic
ordering of the chambers of the type \(A\) arrangement a known or natural object
in Weyl-group theory?

Best regards,

Yehonathan

\end{document}

